% Generated draft: requires mathematical and source review.
\documentclass[11pt]{article}
\usepackage[margin=25mm]{geometry}
\usepackage{fontspec}
\setmainfont{Latin Modern Roman}
\usepackage{amsmath,amssymb,mathtools,unicode-math}
\setmathfont{Latin Modern Math}
\usepackage{enumitem,needspace,xurl,hyperref}
\setlength{\parindent}{0pt}
\setlength{\parskip}{5pt}
\setlength{\emergencystretch}{3em}
\newcommand{\N}{\mathbb{N}}
\newcommand{\Z}{\mathbb{Z}}
\newcommand{\Q}{\mathbb{Q}}
\newcommand{\R}{\mathbb{R}}
\newcommand{\C}{\mathbb{C}}
\newcommand{\F}{\mathbb{F}}
\newcommand{\A}{\mathbb{A}}
\newcommand{\PP}{\mathbb{P}}
\newcommand{\E}{\mathbb{E}}
\newcommand{\eps}{\varepsilon}
\newcommand{\dd}{\,\mathrm{d}}
\begin{document}
\section*{UnsolvedMath ID: 30003546}
\textbf{Source problem code:} catalogue entry
\subsection{Definitions and mathematical statement}
The supplied question is: ``For a prescribed rod structure in the higher-dimensional stationary vacuum construction, can the conical singularities along the bounded axis rods be removed by adjusting the parameters?''

A precise formulation requires a specified construction and its admissible parameter space, neither of which is supplied. The following expresses the balance question conditionally, without asserting that every prescribed rod structure is realizable.

Let \(\mathcal R\) be a prescribed rod structure in the intended construction. A rod structure describes the axis and horizon intervals in the relevant orbit space, together with their associated symmetry data. Its precise meaning, including which geometric data are fixed, must be supplied by that construction. Let \(P_{\mathcal R}\) denote its admissible parameter space and let \(g_\lambda\), \(\lambda\in P_{\mathcal R}\), be its stationary vacuum metrics realizing \(\mathcal R\). Under the usual zero-cosmological-constant convention, the vacuum equation is
\[
\operatorname{Ric}(g_\lambda)=0.
\]
The spacetime dimension, symmetry assumptions, boundary conditions, and admissibility requirements are unspecified in the record.

Write \(I_1,\ldots,I_m\) for the bounded axis rods. Under the usual interpretation that a primitive circle action collapses along each such rod, fix the actual angular identifications of the construction. For \(x\) in the interior of \(I_j\), let \(C_{\lambda,j}(\rho,x)\) be the circumference of the collapsing circle at small proper radial distance \(\rho>0\) from the axis, measured in a spacelike transverse section. Define its transverse cone angle, when this limit exists, by
\[
\alpha_j(\lambda,x)
 =\lim_{\rho\downarrow0}
 \frac{C_{\lambda,j}(\rho,x)}{\rho}.
\]
The conical defect vanishes precisely when this angle is \(2\pi\). This condition concerns conical regularity in the rod interior; it does not by itself establish smoothness at rod endpoints, absence of other singularities, or global regularity.

\textbf{Conditional mathematical target.} For each prescribed rod structure \(\mathcal R\) allowed by the intended construction, determine whether
\[
P_{\mathcal R}^{\mathrm{bal}}
 :=\left\{\lambda\in P_{\mathcal R}:
 \alpha_j(\lambda,x)=2\pi
 \text{ for every }j\in\{1,\ldots,m\}
 \text{ and every }x\in\operatorname{int}(I_j)
 \right\}
\]
is nonempty. Equivalently, determine when one admissible parameter choice simultaneously removes the conical defects along all bounded axis rods while preserving the prescribed rod structure and the construction's other requirements.

\textbf{Unresolved scope.} The record does not determine whether the intended question concerns one particular rod structure, all admissible rod structures, or a general criterion for nonemptiness of \(P_{\mathcal R}^{\mathrm{bal}}\). It also does not specify whether angular periods may vary among the adjustable parameters. Consequently, the conditional formulation above cannot be specialized to a uniquely specified existence problem from the supplied evidence alone.
\subsection{Short English statement}
For a prescribed rod structure, ask whether the allowed parameters can be chosen so that every bounded axis rod has transverse cone angle \(2\pi\) simultaneously. The construction, allowed parameters, angular identifications, dimension, and boundary conditions must be specified before this becomes a definite existence problem. The supplied dated curation labels the problem partially solved; this draft does not independently verify that assessment.
\subsection{Sources}
\begin{itemize}
\item Local UnsolvedMath registry, supplied complete record: permanent ID \(30003546\), source problem code \texttt{catalogue entry}, title ``Removal of Conical Singularities in Higher-Dimensional Vacuum Metrics.'' The record belongs to the Oberwolfach Reports Open Problems collection, gives proposed year 2017, and names no proposer. Its statement is marked \texttt{exact}; the supplied verification describes it as clear, self-contained, publication-ready, and retained byte-for-byte. These are supplied registry assertions.
\item Supplied report attribution: \emph{Analysis, Geometry and Topology of Positive Scalar Curvature Metrics} (2018), Oberwolfach Reports, DOI 10.4171/owr/2017/36, pp.\ 2223--2298. \url{https://doi.org/10.4171/owr/2017/36}. The background describes Oberwolfach Reports as documenting mathematical workshops and their open-problem discussions. No problem-specific page or speaker attribution is supplied, and the connection between this report citation and the stated vacuum-metric question requires review.
\item Supplied literature reference: Lucietti--Tomlinson, ``stationary vacuum black-hole moduli.'' \url{https://arxiv.org/abs/2008.12761}.
\item Supplied literature reference: Lucietti--Tomlinson, ``vacuum black lens obstruction.'' \url{https://doi.org/10.1007/JHEP02(2021)005}.
\item Literature assessment checked at \texttt{2026-08-21}: the supplied curation labels the problem partially solved and reports that integrability and rod-structure methods yield regular higher-dimensional stationary vacuum black holes in many cases, that conical singularities obstruct other rod structures, and that a general parameter-removal criterion is not known. This is supplied dated curation derived from the supplied OWR curation, not an independently verified literature assessment. The record explicitly requests independent verification of time-sensitive status and cited claims.
\end{itemize}
\end{document}
